\section{How the Repository Is Organised}

\begin{verbatim}
distributetr9/
  control/                FastAPI control plane
    app/
      main.py             app startup, background workers
      config.py           every setting, read from environment variables
      models.py           SQLAlchemy tables
      db.py               engine, sessions, migrations on startup
      scheduler.py        claim, renew, complete, lease expiry
      envbuild.py         custom environment builds
      agent_bundle.py     signed agent updates, install package
      alerts.py ops.py    operations: alerts, garbage collection
      routers/            HTTP endpoints: portal.py, agent.py, auth.py
    migrations/           Alembic schema migrations
  agent/                  runs on each lab PC (Python standard library only)
    agent.py              main loop: register, heartbeat, claim, run
    executor.py           docker run / local python, log pumping
    classify.py           GPU detection and class mapping
    updater.py            verify and install signed updates
    launcher.sh run.sh    process supervision and rollback
    classes.json          GPU class catalog
  portal/                 Next.js web UI
  harness/                fake agents for tests and load testing
  tests/                  pytest suite (runs against real Postgres + MinIO)
  sim/                    simulated five-PC lab + failure exam
  deploy/                 production compose, Caddy, backup, restore
  scripts/                dev launcher, lab PC installer, smoke test
  supervisor/             Windows boot task and one-command bootstrap
  spec/                   the design specification (14 modules, 12 phases)
  docs/                   runbooks, registry notes, this guide
\end{verbatim}

\subsection{Why this split}

Each directory matches one \emph{deployable thing} or one \emph{audience}:

\begin{itemize}
\item \texttt{control/} and \texttt{portal/} run on the server.
\item \texttt{agent/}, \texttt{supervisor/} and the installer run on lab PCs.
  The agent deliberately uses only the Python standard library, so installing
  it means copying files, and updating it means copying newer files. There is
  no \texttt{pip install} on sixty machines that might each fail differently.
\item \texttt{harness/}, \texttt{tests/} and \texttt{sim/} exist only to
  prove the other parts work.
\item \texttt{spec/} was written before the code. It froze the scope of version
  1 and split the work into phases, each ending with a \emph{gate}: a concrete
  demonstration that must pass before the next phase starts.
\end{itemize}

\subsection{Configuration: one settings object}

Every tunable value lives in \texttt{control/app/config.py} as a field on a
\texttt{pydantic-settings} class. Pydantic reads each field from an environment
variable of the same name (case-insensitive), falling back to a
\texttt{.env} file and then to the default:

\begin{lstlisting}[language=Python]
class Settings(BaseSettings):
    model_config = SettingsConfigDict(env_file=".env", extra="ignore")

    lease_ttl_seconds: int = 30
    claim_poll_seconds: int = 25
    agent_rollout_ring: int = 2

settings = Settings()   # LEASE_TTL_SECONDS=10 in the environment wins
\end{lstlisting}

The benefit is that tests, the simulated lab and production run the same code
with different numbers. The simulated lab, for example, sets a 10 second lease
so that failure scenarios resolve in seconds rather than minutes.

\paragraph{Other use cases.} This is the ``config in the environment'' rule of
the Twelve-Factor App methodology, used by most containerised services. The
same pattern appears in Django (\texttt{django-environ}), Spring Boot
(\texttt{application.properties} overridden by environment variables) and
Kubernetes ConfigMaps.

\subsection{The development workflow}

\begin{lstlisting}[language=bash]
make dev        # Postgres, MinIO, registry, API :8080, portal :3000, agents
make test       # pytest against a throwaway lab_test database
make lint       # ruff: syntax errors, undefined names, unused imports
make sim        # five simulated lab PCs in Docker
make sim-exam   # break them in thirteen different ways
\end{lstlisting}

A \texttt{Makefile} here is only a list of named shortcuts; it does not build
anything. It helps newcomers find the important commands without reading the
README.
